\documentclass[11pt]{article}
\usepackage[margin=0.85in]{geometry}
\usepackage{hyperref}
\usepackage{enumitem}
\usepackage{titlesec}
\usepackage{parskip}
\usepackage[T1]{fontenc}

\hypersetup{colorlinks=true, urlcolor=blue, linkcolor=blue}

\titleformat{\section}{\Large\normalfont}{}{0em}{}[\titlerule]
\titlespacing*{\section}{0pt}{12pt}{6pt}

\setlist[itemize]{leftmargin=*, itemsep=2pt, topsep=2pt}

\pagenumbering{gobble}

\newcommand{\entry}[4]{%
  \noindent\textbf{#1}\hfill\textit{#2}\\
  \textit{#3}\vspace{-2pt}
  #4
}

\begin{document}

\begin{center}
{\Huge Alexander Pino}\\[4pt]
\href{mailto:alex.pino.salig@gmail.com}{alex.pino.salig@gmail.com} \quad | \quad \href{https://linkedin.com/in/alexander-j-pino}{linkedin.com/in/alexander-j-pino} \quad | \quad \href{https://github.com/a-pino}{github.com/a-pino}
\end{center}

\section*{Experience}

\entry{Software Engineer}{Jun 2025 -- Present}{Accenture Industry X}{
\begin{itemize}
\item Owns the maintenance and deployment of \textbf{AIOHTF}, an internal hardware-testing platform extending Google's OpenHTF with an asynchronous FastAPI layer and parallel multi-plug configuration; integrated open-source YOLO and Qwen-VL models (PyTorch) for automated validation of embedded software UIs, with annotation pipelines managed in Label Studio. \textit{(Client: major U.S. manufacturer of industrial water treatment and sanitation equipment.)}
\item Led a five-engineer team prototyping \textbf{AIX Agent}, an agentic platform with a FastAPI backend and React frontend that enables an AI agent to autonomously manage physical industrial sites on Azure---querying embedded sensor data from SQL Server, performing root cause analysis, enforcing SOPs for issue triage and shift handoff, and generating compliance and shift-handoff reports.
\end{itemize}
}

\entry{Research Intern}{Jul -- Sep 2024}{AIKOCorp}{
\begin{itemize}
\item Architected and implemented an automated privilege-escalation system combining a graph neural network (GNN) with a reinforcement learning (RL) environment (PyTorch, HuggingFace) to autonomously obtain administrator access within simulated Active Directory networks.
\end{itemize}
}

\entry{Research Fellow}{Feb -- Jun 2024}{WhiteBox Research}{
\begin{itemize}
\item Designed a benchmarking framework (PyTorch, HuggingFace) evaluating frontier and open-source LLMs on their self-awareness of safety risks when answering queries related to Chemical, Biological, Radiological, and Nuclear (CBRN) topics; conducted in collaboration with interpretability research institute Apart Research.
\end{itemize}
}

\entry{Research Intern}{Jul 2023}{FrontierLab@OsakaU}{
\begin{itemize}
\item Built a computer vision pipeline (PyTorch, HuggingFace) for multi-species microorganism identification in petri dish imagery, benchmarking generalist models (YOLO, ViT) against microscopy-specialized models (StarDist, Cellpose). Supervised by Hideo Matsuda, PhD and Shigeto Seno, PhD; presented research poster on July 31, 2023. \textit{(Synthetic Ecosystem Project, Institute of Science and Technology.)}
\end{itemize}
}

\section*{Education}

\entry{BS Applied Mathematics --- Specialization in Data Science}{Jun 2024}{Ateneo de Manila University}{
\begin{itemize}
\item \textbf{Thesis:} \textit{Using Topology to Extract Insights from UAAP Basketball Data} (with B.L. Ang and D.L. Rosario). Applied Vietoris-Rips filtration and the Mapper algorithm to UAAP player statistics, discovering four novel player subgroups (paint dominators, stretch forwards, defensive rebounders, scrappers) and identifying two underrepresented archetypes within the league.
\end{itemize}
}

\end{document}
